\pdfbookmark[0]{Magnetisches Observatorium. 1834}{magnetisches_observatorium_1834}
\begin{center}
\hrule\vspace{2mm}
Göttingische gelehrte Anzeigen. 1834. 9. August.
\vspace{2mm}\hrule\vspace{5mm}
\end{center}
Wir verdanken dem Wohlwollen unserer Regierung ein neues, einem wichtigen Teil der Naturwissenschaften gewidmetes Institut, ein eigens
\begin{center}\itshape für die magnetischen Beobachtungen und Messungen\end{center}
\noindent errichtetes Observatorium.

Obwohl der Bau bereits im vergangenen Herbst und die Einrichtung im Inneren seit Anfang dieses Jahres so weit abgeschlossen war, dass schon seit den ersten Monaten tägliche Beobachtungen möglich waren, haben wir bisher gezögert, in diesen Blättern darüber zu berichten. Wir wollten erst einige Beobachtungsergebnisse mit dem Bericht verbinden können. Die nach neuen Prinzipien gebauten magnetischen Apparate, die 1832 in der hiesigen Sternwarte aufgestellt wurden, haben wir bereits früher in diesen Blättern [Anzeige der \textit{Intensitas v.\,m.}] ausführlich beschrieben, und die damit
\SourceLeaf{526}erreichbare Genauigkeit geht aus der damaligen Darstellung hinreichend hervor. Um diese Genauigkeit jedoch \textit{voll} zu erreichen, war eine Ausführung in größerem Maßstab nötig. Um die Ergebnisse völlig frei von fremden Einflüssen zu halten, war ein eigenes eisenfreies Gebäude unerlässlich.

Das magnetische Observatorium wurde auf einem freien Platz etwa hundert Schritte westlich der Sternwarte errichtet. Es ist ein genau ausgerichtetes längliches Rechteck von 32 Pariser Fuß Länge und 15 Fuß Breite, mit zwei Vorsprüngen an den Längsseiten. Der westliche Vorsprung bildet den Eingang und dient zugleich bei bestimmten Beobachtungen als Erweiterung des Hauptsaals. Der östliche Vorsprung ist vom Hauptsaal ganz getrennt und dient dem Nachtwächter der Sternwarte als Aufenthaltsraum. Im ganzen Gebäude besteht ausnahmslos alles aus Kupfer, wofür sonst Eisen verwendet wird, etwa Schlösser, Türangeln, Fensterbeschläge und Nägel. Soweit möglich wurde dafür gesorgt, jegliche Zugluft fernzuhalten. Der Saal ist etwas über 10 Fuß hoch.

Der magnetische Apparat stimmt im Wesentlichen mit den oben erwähnten überein. Wir beschränken uns daher darauf, die Unterschiede anzugeben. Der Magnetstab besteht aus Uslarer Gussstahl, der sich ausgezeichnet für magnetische Versuche eignet. Von Zeit zu Zeit werden verschiedene Stäbe eingesetzt, die alle nahezu gleich groß sind, nämlich 610 Millimeter lang, 37 breit und 10 dick. Sie wiegen etwa vier Pfund. Der Spiegel ist 75 Millimeter breit und 50 hoch. Der Stab hängt von der Mitte der Saaldecke an einem 200-fachen, 7 Fuß langen, ungedrehten Seidenfaden. Der Torsionskreis sitzt jedoch anders als früher am unteren Ende des Fadens statt am oberen und ist mit dem Schiffchen, das den Stab trägt, drehbar verbunden. Wie Hofrat \textsc{Gauss} in seiner Abhandlung (\textit{Intensitas vis magneticae terrestris}, Art.\ 9) bereits bemerkt hat, haben seidene Aufhängungsfäden gegenüber metallenen den großen Vorteil, dass ihre Torsionskraft sehr klein ist. Beim gegenwärtigen Tragfaden beträgt sie nur ein Neunhundertstel der horizontalen Richtkraft des Magnetstabs, während die Torsionskraft eines Metallfadens mit gleicher Tragfähigkeit etwa zehnmal so groß wäre. Seidenfäden haben dagegen, besonders wenn ihre Tragfähigkeit das daran hängende Gewicht nicht weit übersteigt, den Nachteil, dass sie sich in den ersten Wochen oder bei erheblich verstärkter Belastung beträchtlich verlängern. Diesem Nachteil hilft hier der sinnreiche, von Professor \textsc{Weber} angege\SourceLeaf{527}bene Aufhängungsapparat an der Decke ab. Mit ihm lässt sich der Faden leicht so weit wie nötig wieder aufwinden, ohne seinen Platz zu verändern. Zugleich lässt sich der Apparat ebenso leicht an der Decke verschieben, wenn die Änderung der magnetischen Deklination dies mit der Zeit nötig macht. Der Theodolit steht bisher auf einem sehr solide gearbeiteten Holzstativ über einem eigenen steinernen Fundament. Von seinem Standort aus ist durch das Nordfenster einer der Stadttürme zu sehen, dessen Azimut mit größter Genauigkeit bestimmt wurde. Zur Kontrolle der unveränderten Stellung des Theodoliten dient allein ein feiner senkrechter Strich an der gegenüberliegenden Nordwand. Für den gewöhnlichen Gebrauch dient eine 4 Fuß lange, in Millimeter geteilte Skala. Bei manchen Beobachtungen wird sie durch eine zwei Meter lange ersetzt. Der Wert eines Skalenteils beträgt $21''3$. Bei nächtlichen Beobachtungen wurde die Skala bisher mit kräftigen Wachskerzen beleuchtet. Künftig werden dazu \textsc{Argand}-Lampen verwendet.

Eine der Hauptanwendungen des Apparats ist die genaue Bestimmung der magnetischen Deklination und ihrer Änderung zu verschiedenen Tagesstunden, in verschiedenen Monaten und Jahren. Jeden Tag wird die Ablesung zweimal zu festen Stunden aufgezeichnet. Gewählt wurden 8 Uhr vormittags und 1 Uhr nachmittags. Bei regelmäßigem Verlauf der täglichen Schwankungen fallen die kleinste und die größte Deklination zumindest in den ersten Monaten des Jahres ungefähr mit diesen Zeiten zusammen. Aus wichtigen Gründen, die hier nicht weiter ausgeführt werden können, musste man unbedingt der Aufzeichnung zu stets derselben Uhrzeit den Vorzug geben, statt jedes Mal das Minimum und das Maximum abzuwarten. Mit diesen Aufzeichnungen wurde zwar schon am 1. Januar begonnen. Zunächst verwendete man jedoch einen schwächeren Aufhängungsfaden, der sich nach und nach verlängerte und daher öfter aufgewunden werden musste. Dabei traten nicht unerhebliche Änderungen des Nullpunkts der Torsion ein, die man anfangs zu wenig beachtete. Man hat es deshalb vorgezogen, die ersten zweieinhalb Monate auszuschließen. Die seither erhaltenen Mittelwerte für die westliche Deklination der Magnetnadel waren folgende.
\SourceLeaf{528}
\begin{center}
\begin{tabular}{l|rrr|rrr}
 & \multicolumn{3}{c|}{8 Uhr vorm.} & \multicolumn{3}{c}{1 Uhr nachm.}\\\hline
März, zweite Hälfte & $18^0$ & $38'$ & $16''0$ & $18^0$ & $46'$ & $40''4$\\
April & & $36$ & $6.9$ & & $47$ & $3.8$\\
Mai & & $36$ & $28.2$ & & $47$ & $15.4$\\
Juni & & $37$ & $40.7$ & & $47$ & $59.5$\\
Juli & & $37$ & $57.5$ & & $48$ & $19.0$
\end{tabular}
\end{center}

Ferner werden an bestimmten Tagen im Jahr 44 Stunden lang ohne Unterbrechung in kurzen Zeitabständen die Änderungen der Deklination beobachtet. Dafür wählte man dieselben Tage, die Herr von \textsc{Humboldt} bereits vor mehreren Jahren festgelegt hatte. An diesen Tagen werden nach Absprache schon an vielen, zum Teil sehr weit entfernten Orten ähnliche Aufzeichnungen mit \textsc{Gambey}-Apparaten gemacht. Bis jetzt wurden diese Beobachtungen hier dreimal angestellt, nämlich am 20. und 21. März, am 4. und 5. Mai und am 21. und 22. Juni. Daran nahmen außer Hofrat \textsc{Gauss} die Herren Professor \textsc{Weber}, Professor \textsc{Ulrich}, Dr.\ \textsc{Weber}, Dr.\ \textsc{Goldschmidt}, Dr.\ \textsc{Listing}, \textsc{Sartorius}, \textsc{Deahna} und \textsc{Wilh.\ Gauss} teil. Die Beobachtungen sollen zum einen den regelmäßigen Verlauf nach und nach immer vollständiger erkennen lassen. Zum anderen soll der Vergleich gleichzeitiger Beobachtungen an verschiedenen Orten zeigen, was es mit den außerordentlichen Anomalien auf sich hat, die so häufig dazwischenkommen und zuweilen, besonders bei Nordlichtern, außerordentlich groß sind. Hier wurde im März von 20 zu 20 Minuten aufgezeichnet, zum Teil in halb so großen Abständen, im Mai von 10 zu 10 Minuten und zum Teil doppelt so häufig, im Juni durchgehend von 5 zu 5 Minuten. Ein paar auffallend große Anomalien wurden hier in der Nacht vom 20. zum 21. März bemerkt. In den Nächten vom 4. und 5. Mai traten sehr erhebliche und zahlreiche auf. Am 21. Juni waren einige zwar nicht groß, traten aber deutlich hervor, während der Verlauf am ganzen 22. Juni überaus regelmäßig war. Von den entsprechenden Beobachtungen, die, wie erwähnt, auf die Anregung von Herrn von \textsc{Humboldt} zurückgehen, sind uns bisher nur die Berliner vom 20. und 21. März bekannt geworden. Sie wurden jedoch nur stündlich aufgezeichnet und konnten daher keine besonderen Ergebnisse liefern, obwohl sich in ihnen die in Göttingen bemerkten und verfolgten Anomalien andeuteten. Herr \textsc{Sartorius} verwendete dagegen einen zwar kleineren, aber nach denselben Prinzipien wie
\SourceLeaf{529}der hiesige gebauten Apparat, um am 4. und 5. Mai auf einem Gut in Bayern, einige Meilen südlich von Meiningen, eine sehr vollständige Reihe entsprechender Beobachtungen anzustellen. Daraus ergibt sich eine wahrhaft bewundernswerte Übereinstimmung mit den hier beobachteten großen Anomalien nach Zeit, Größe und Wechsel. In den grafischen Darstellungen möchte man die eine beinahe für eine Kopie der anderen halten, mit all den bizarren Formen, die durch jene Anomalien hervortreten. Ein ebenso schönes Ergebnis zeigte sich am 21. und 22. Juni. In Berlin stellte Professor \textsc{Encke} mit Unterstützung der Herren \textsc{Poggendorff}, \textsc{Mädler} und \textsc{Wolfers} zum ersten Mal entsprechende Beobachtungen mit einem dem hiesigen ähnlichen, wenn auch kleineren Apparat an. Auch dort gab es keine anderen Anomalien als die hier beobachteten, aber diese waren fast genau wiedergegeben. Genau dasselbe zeigten die Beobachtungen, die Herr \textsc{Sartorius} dieses Mal in Frankfurt am Main machte. Diese Ergebnisse können schon als schöne Frucht der verabredeten Beobachtungen gelten. Aus ihnen geht aufs Klarste hervor, dass kleinere und größere Anomalien der Magnetnadel, die zuweilen in recht kurzen Abständen wechseln, keine örtlichen, sondern kräftige, weithin wirkende Ursachen haben müssen. Das war bereits früher bei sehr großen Unregelmäßigkeiten im Zusammenhang mit Nordlichtern bemerkt worden. Wenn sich die Teilnahme an diesen verabredeten Beobachtungen mit den ebenso genauen wie bequemen Apparaten künftig immer weiter ausbreitet, wofür schon jetzt die besten Aussichten bestehen, werden wir gewiss umfassenden Aufschluss über diese höchst bemerkenswerten und rätselhaften Erscheinungen erhalten.

Solche Beobachtungen werden hier im Übrigen auch außerhalb der festgelegten Zeiten häufig angestellt. Dabei sind zuweilen ganz auffallende Anomalien vorgekommen. So nahm etwa am 14. Januar abends zwischen 8 und 9 Uhr die Deklination innerhalb einer einzigen Viertelstunde mit größter Regelmäßigkeit um 13 Bogenminuten ab und kehrte dann allmählich auf ihren vorigen Stand zurück. Solche Beobachtungen können jedoch keine weiteren Ergebnisse liefern, da ohne Absprache nur äußerst selten entsprechende Beobachtungen zu erwarten sind.

Von Zeit zu Zeit wird im hiesigen magnetischen Observatorium auch die Bestimmung der absoluten Intensität des Erdmagnetismus wiederholt werden. Um diese Messung mit größter Genauigkeit durchzuführen, mussten erst verschie\SourceLeaf{530}dene Vorkehrungen getroffen werden. Deshalb konnte sie zum ersten Mal erst im Juli vorgenommen werden. Drei Bestimmungen mit verschiedenen Stäben ergaben
\begin{center}
\begin{tabular}{rllllllr}
17. Juli & . & . & . & . & . & & $1.7743$\\
20. \enspace „ & . & . & . & . & . & & $1.7740$\\
21. \enspace „ & . & . & . & . & . & & $1.7761$
\end{tabular}
\end{center}
als Wert der horizontalen Kraft. Dabei liegen wie bei den früheren Bestimmungen mit kleineren Stäben die Sekunde als Zeiteinheit, der Millimeter und das Milligramm zugrunde. Man wird mit Freude bemerken, wie wenig die früheren von den gegenwärtigen Ergebnissen abweichen.

Wie beim früheren, in der Sternwarte aufgestellten Apparat hat man nun auch den gegenwärtigen Apparat im magnetischen Observatorium mit Vorrichtungen für elektromagnetische Versuche und Messungen verbunden. Der aufgehängte Magnetstab ist von einem Multiplikator mit 200 Windungen umgeben, dessen Bauart die Verwendung von nicht umsponnenem Draht erlaubte. Die Drahtlänge beträgt 1100 Fuß. Mit einem sehr einfach gebauten Kommutator kann der Beobachter jeden Augenblick die Richtung des galvanischen Stroms umkehren oder den Strom ganz unterbrechen, ohne das Auge vom Fernrohr zu nehmen.

Wir können dabei eine großartige und bisher in ihrer Art einzige Anlage nicht unerwähnt lassen, die mit den beschriebenen Einrichtungen eng verbunden ist und die wir unserem Professor \textsc{Weber} verdanken. Er hatte schon im vergangenen Jahr vom physikalischen Kabinett aus über die Häuser der Stadt hinweg bis zur Sternwarte eine doppelte Drahtverbindung geführt. Sie ist jetzt von der Sternwarte bis zum magnetischen Observatorium verlängert. So entsteht eine große galvanische Kette, in der der galvanische Strom, die Multiplikatoren an beiden Endpunkten mitgerechnet, fast neuntausend Fuß Draht durchlaufen muss. Der Draht der Kette ist größtenteils Kupferdraht der Handelsnummer 3, von dem ein Meter acht Gramm wiegt. Der Draht des Multiplikators im magnetischen Observatorium ist versilberter Kupferdraht Nr.\ 14, von dem 2,6 Meter auf ein Gramm kommen. Diese Anlage eignet sich ausgezeichnet dazu, Gelegenheit zu einer Fülle höchst interessanter Versuche zu geben. Man be\SourceLeaf{531}merkt nicht ohne Bewunderung, wie ein einziges am anderen Ende eingesetztes Plattenpaar den Magnetstab augenblicklich in Bewegung setzt, bis zu einem Ausschlag von weit über tausend Skalenteilen. Noch auffallender findet man zumindest anfangs, dass ein sehr kleines Plattenpaar, zum Beispiel von einem einzigen Zoll Durchmesser, mit bloßem Brunnenwasser oder sogar destilliertem Wasser eine nicht viel kleinere Wirkung hervorbringt als ein sehr großes Plattenpaar mit starker Säure. Bei näherem Nachdenken ist dies jedoch ganz in der Ordnung und bestätigt nur aufs Neue die schöne, zuerst von \textsc{Ohm} aufgestellte Theorie. Erhöht man dagegen die Anzahl der Plattenpaare, so wächst die Wirkung, und zwar fast proportional zu dieser Anzahl. Die Leichtigkeit und Sicherheit, mit der man durch den Kommutator die Stromrichtung und die davon abhängige Bewegung der Nadel beherrscht, hatte schon im vergangenen Jahr Versuche zur telegrafischen Signalübermittlung veranlasst. Diese gelangen auch mit ganzen Wörtern und kurzen Wendungen aufs Vollkommenste. Es besteht kein Zweifel, dass sich auf ähnliche Weise eine unmittelbare telegrafische Verbindung zwischen zwei Orten einrichten ließe, die eine beträchtliche Zahl von Meilen voneinander entfernt sind. Dies ist aber selbstverständlich nicht der Ort, Gedanken darüber weiter auszuführen.

\bigskip
